\documentclass[11pt,a4paper]{article}
\usepackage[margin=1in]{geometry}
\usepackage{hyperref}
\usepackage{enumitem}
\usepackage{graphicx}
\usepackage{xcolor}
\usepackage{booktabs}
\usepackage{array}

% -----------------------------------------------------
% bvraghav-tiet-ucs503p-article.sty
% -----------------------------------------------------

% Variable: Lab Instructor
\makeatletter \newcommand{\BvrSetLabInstructor}[1]{%
  \gdef\Bvr@LabInstructor{#1}}
% default
\newcommand{\Bvr@LabInstructor}{Dr. Raghav %
  B. Venkataramaiyer}
% public accessor
\newcommand{\BvrLabInstructorValue}{\Bvr@LabInstructor}
\makeatother

% Add subtitle
\usepackage{titling}
% -- Set title bold
\pretitle{\begin{center}\bfseries\fontsize{18bp}{18bp}\selectfont}
\makeatletter
\newcommand{\BvrSubtitle}[1]{%
  \posttitle{%
    \par\end{center}
    \begin{center}\large#1\end{center}
    \vskip 0.5em}%
}
\makeatother

% Hints in the section title(s)
\newcommand{\BvrHint}[1]{%
  {\normalfont\normalsize\itshape\hfill #1}}

% -----------------------------------------------------

\title{API Gateway: A Self-Service Rate-Limiting Layer for Web APIs}

\BvrSetLabInstructor{Dr. Paramveer Sidhu}

\BvrSubtitle{%
  Project Proposal for UCS503P  \\
  Submitted to: \BvrLabInstructorValue \\ \bfseries \\[0.5em]
  Thapar Institute of Engineering and Technology%
}

\author{
  Sarthak Jagota - \texttt{1024030753} \\[0.3em]
  Palak Mahajan - \texttt{1024030755} \\[0.3em]
  Aayush Kumar Singh - \texttt{1024030758}
}

\date{\today}

\begin{document}
\maketitle

\section{Higher-order goal%
  \BvrHint{\ldots secondary framing}}
This project aims to improve the \emph{resilience and cost-predictability of
networked services} by making traffic control available to developers who
cannot afford, or do not wish to operate, enterprise API-management
infrastructure. While traffic shaping is a broad concern spanning cloud
economics and denial-of-service resistance, the immediate engineering
objective is narrower and concrete: to translate \emph{fair, enforceable
request budgeting} into a deployable software layer that any API owner can
place in front of an existing service without modifying that service.

\section{Time-to-value %
\BvrHint{\ldots primary heuristic}}
\begin{itemize}[leftmargin=0.7cm]
    \item The smallest useful product is a proxy that forwards a request and
      returns HTTP~429 when a per-client budget is exhausted. This is
      demonstrable in the first iteration and is independently valuable even
      with no dashboard, no persistence and a single algorithm.
    \item We will prioritise fast feedback over completeness by using weekly
      iterations with CI-backed automation from the outset, so that each
      increment is releasable to staging rather than accumulating on a branch.
    \item The work splits into three parallel workstreams --- gateway and
      rate-limiting logic, persistence and authentication, and the analytics
      frontend --- allowing the team to produce measurable increments
      concurrently rather than serially.
    \item Success in iteration one is defined by measurement, not by feature
      count: enforcement must be observably correct under concurrent load
      before any feature is layered on top of it.
\end{itemize}

\section{Problem Statement}
An HTTP API, by default, has no notion of a per-client budget. Every client is
free to issue requests as fast as it can open connections, and the server
discovers the consequences only after they occur. This produces four
recurring failure modes:
\begin{itemize}[leftmargin=0.7cm]
    \item \textbf{Resource exhaustion:} a single misbehaving client---whether
      a retry loop caused by a bug, a scraper, or a deliberate abuser---can
      saturate connection pools, database connections or worker threads, and
      degrade service for every other client.
    \item \textbf{Uncontrolled cost:} where the API is backed by metered
      resources (managed databases, third-party APIs, egress bandwidth),
      unbounded request volume translates directly into unbounded expenditure,
      discovered only at billing time.
    \item \textbf{No fairness guarantee:} without per-client accounting, a
      well-behaved consumer and an abusive one are indistinguishable to the
      server, so capacity is allocated by aggression rather than by policy.
    \item \textbf{High cost of the remedy:} the standard solution---a
      rate-limiting gateway between clients and the origin---exists in mature
      form (Kong, Cloudflare, AWS API Gateway) but is either operationally
      heavy or commercially priced. A developer wishing to protect a single
      API must otherwise build the counting, storage and enforcement logic
      themselves.
\end{itemize}

\textbf{Why this matters now (contextual relevance):} Deployment on shared
platform-as-a-service tiers has made small APIs cheap to publish and
correspondingly easy to overwhelm; the marginal cost of a request has moved
from the developer's fixed server bill to a metered per-request charge. The
absence of a lightweight, self-service traffic-control layer is therefore a
present and practical gap rather than a hypothetical one.

\section{Proposed Solution %
\BvrHint{\ldots Software Proposition}}
\subsection{Overview}
We propose \textbf{API Gateway}, a web-based API protection layer that sits as a
reverse proxy between a client and a target API. Its components are:
\begin{itemize}[leftmargin=0.7cm]
    \item \textbf{Gateway/proxy layer:} intercepts every inbound request,
      evaluates it against the applicable policy, forwards permitted requests
      to the registered origin and rejects the remainder with HTTP~429
      (\texttt{Too Many Requests}).
    \item \textbf{Self-service registration portal:} an authenticated user
      registers a target API URL, sets a request budget (requests per unit
      time), and selects the enforcement algorithm.
    \item \textbf{Distributed counting substrate:} per-client counters held in
      Redis and manipulated through Bucket4j, so that the limit remains
      correct when several gateway instances serve the same API.
    \item \textbf{Analytics dashboard:} a React interface showing allowed
      versus blocked volume, traffic over time, and a per-API breakdown.
    \item \textbf{Request log:} a queryable table of recent requests with
      outcome, timestamp and client identity, forming the evidential basis for
      both the dashboard and our evaluation metrics.
\end{itemize}

\subsection{Core workflow}
\begin{enumerate}[leftmargin=0.7cm]
    \item A developer signs up, authenticates, and registers a target API,
      receiving a gateway endpoint and an API key.
    \item A client issues a request to that gateway endpoint bearing the key.
    \item The gateway resolves the policy for (client, API), and performs an
      atomic consume-or-reject operation against the Redis-backed counter.
    \item If a token is available the request is proxied to the origin and the
      origin's response is relayed verbatim; otherwise the gateway short-
      circuits with HTTP~429 and the applicable \texttt{Retry-After} hint.
    \item In both cases an outcome record is written to PostgreSQL.
    \item The dashboard reads aggregated records over REST and renders them.
\end{enumerate}

\subsection{Operational constraints for a course-project setting}
\begin{itemize}[leftmargin=0.7cm]
    \item The gateway must add only bounded latency to the proxied path; a
      correct limiter that doubles response time is not a usable one.
    \item We avoid dependence on novel research: token bucket and sliding
      window are well-specified, published algorithms, and our contribution is
      their correct distributed implementation and instrumentation, not their
      invention.
    \item The whole system must deploy on free or low-cost managed tiers
      (application host, managed Redis, managed PostgreSQL) so that the pilot
      is reproducible by the evaluator without provisioning cost.
    \item The dashboard must remain usable on low-to-mid range devices through
      responsive design and server-side aggregation rather than shipping raw
      logs to the browser.
\end{itemize}

\section{Solution Approach %
\BvrHint{\ldots Engineering focus}}
This is an engineering course project. Accordingly the emphasis falls on
correctness under concurrency, clean separation of concerns, and measurable
behaviour, rather than on algorithmic novelty.

\subsection{Rate-limiting strategy %
\BvrHint{\ldots non-research}}
Enforcement uses two standard algorithms, both provided by Bucket4j over a
Redis backend:
\begin{itemize}[leftmargin=0.7cm]
    \item \textbf{Token bucket:} a bucket of capacity $C$ refilled at rate $r$;
      each request consumes one token, and a request arriving at an empty
      bucket is rejected. This permits controlled bursts up to $C$ while
      bounding the long-run average at $r$.
    \item \textbf{Sliding window:} the count of requests within a moving
      interval of width $W$ is compared against a threshold $N$, giving
      smoother enforcement at the cost of more state per client.
    \item \textbf{Correctness under distribution:} counter updates are
      performed as atomic server-side operations in Redis, so that two gateway
      instances handling simultaneous requests from the same client cannot
      both observe the same available token. This---not the arithmetic of
      either algorithm---is the substantive engineering problem.
    \item \textbf{Failure policy:} the behaviour when the counter store is
      unreachable is an explicit design decision (fail-open, preserving
      availability, versus fail-closed, preserving protection). We will
      implement one, document the rationale, and test it.
\end{itemize}

\subsection{Web architecture %
\BvrHint{\ldots web-based solution}}
A three-tier architecture with an additional in-memory tier for counters:
\begin{itemize}[leftmargin=0.7cm]
    \item \textbf{Frontend:} React single-page application with Recharts for
      visualisation; communicates with the backend exclusively over REST.
    \item \textbf{Gateway/Backend:} Spring Boot with Spring Cloud Gateway for
      the proxy path; Spring Security with JWT for authentication; Spring Data
      JPA for persistence.
    \item \textbf{Counter store:} Redis, holding only ephemeral per-client
      counter state.
    \item \textbf{Database:} PostgreSQL, holding users, registered APIs,
      policies and request logs.
\end{itemize}

\noindent The request path is summarised below.

\begin{center}
\begin{minipage}{0.92\textwidth}
\ttfamily\small
Client\\
\quad$\downarrow$\\
API Gateway (Spring Cloud Gateway)\\
\quad$\downarrow$\ \ consume-or-reject\\
Bucket4j + Redis (distributed counter)\\
\quad$\downarrow$\ \ if allowed\\
Target API \ $\longrightarrow$ \ response relayed to client\\
\quad$\downarrow$\ \ always\\
Request log written \ $\longrightarrow$ \ PostgreSQL (Spring Data JPA)\\
\quad$\downarrow$\\
React Dashboard (REST) \ $\longrightarrow$ \ charts and logs
\end{minipage}
\end{center}

\subsection{Technology selection}
\begin{center}
\begin{tabular}{@{}>{\raggedright\arraybackslash}p{0.34\textwidth}>{\raggedright\arraybackslash}p{0.56\textwidth}@{}}
\toprule
\textbf{Layer} & \textbf{Technology} \\
\midrule
Gateway / Backend        & Spring Boot, Spring Cloud Gateway \\
Rate limiting            & Bucket4j \\
Counter store            & Redis \\
Primary database         & PostgreSQL (Spring Data JPA) \\
Authentication           & Spring Security + JWT \\
Frontend                 & React + Recharts \\
API testing              & Postman; JUnit for automated tests \\
Deployment               & Render/Railway (backend, Redis, PostgreSQL); Vercel (frontend) \\
\bottomrule
\end{tabular}
\end{center}

\subsection{CI/CD and fast delivery enabler}
CI/CD will be used to:
\begin{itemize}[leftmargin=0.7cm]
    \item build the backend and run unit and integration tests on every change;
    \item run the load-and-enforcement test suite, so that a regression in
      limiting accuracy fails the build rather than reaching staging;
    \item produce repeatable deployment artifacts (container images) for
      staging;
    \item enable safe and frequent release of incremental features.
\end{itemize}

\section{Evaluation Criterion %
\BvrHint{\ldots measurable and attributable}}
Success is defined by instrumented measurement on the deployed system, not by
feature presence.

\subsection{Primary evaluation metric}
\textbf{Enforcement accuracy under concurrent load (EA):} for a configured
limit of $N$ requests per window, the number of requests actually admitted to
the origin when $M \gg N$ concurrent requests are issued from a single client
identity across all running gateway instances.
\begin{itemize}[leftmargin=0.7cm]
    \item Target: admitted count within $\pm 2\%$ of $N$, with two gateway
      instances active, over a driver issuing at least $10N$ requests.
    \item Attribution: computed by comparing the load driver's own tally of
      HTTP~200 versus HTTP~429 responses against the count of forwarded
      requests recorded at the origin. Over-admission indicates a race in the
      counter path; under-admission indicates an over-strict window boundary.
      Both are attributable to the limiter, not to the network.
\end{itemize}

\subsection{Secondary evaluation metrics}
\begin{itemize}[leftmargin=0.7cm]
    \item \textbf{Added latency (p50/p95):} difference in response time between
      a request routed through the API Gateway and the same request issued directly
      to the origin. Target: p95 overhead $\le$ 50\,ms on the pilot deployment.
    \item \textbf{Sustained throughput:} requests per second the gateway
      sustains before p95 latency degrades beyond the above threshold,
      reported per instance count to evidence horizontal scaling.
    \item \textbf{Log completeness:} proportion of driver-issued requests
      appearing in the PostgreSQL log with the correct outcome. Target:
      $\ge$ 99.9\%, since the dashboard's credibility rests on this.
    \item \textbf{Dashboard freshness:} lag between a request occurring and its
      appearance in dashboard aggregates.
    \item \textbf{Reliability:} availability of the staging gateway for proxy
      and login operations during the pilot window (target $\ge$ 99\% uptime).
\end{itemize}

\subsection{Pilot validation plan %
\BvrHint{\ldots high-level}}
\begin{itemize}[leftmargin=0.7cm]
    \item Register a small set of representative target APIs (a deliberately
      slow endpoint, a fast endpoint, and a public third-party API) to
      exercise differing latency profiles.
    \item Drive load using a scripted client at controlled concurrency levels,
      repeating each run with one and with two gateway instances to
      demonstrate that distributed counting holds.
    \item Recruit 5--10 developer users to register their own API and report
      clarity of setup and of the dashboard on a 1--5 scale, providing a
      usability counterweight to the purely quantitative metrics.
\end{itemize}

\section{Scalability %
\BvrHint{\ldots with foresight}}
The proposition should scale in both theoretical and practical senses.

\begin{enumerate}[leftmargin=0.7cm]
    \item \textbf{Scalable (theoretically):} the design admits growth in
      request volume, registered APIs and concurrent dashboard users by:
    \begin{itemize}[leftmargin=0.7cm]
        \item keeping the gateway \emph{stateless}, with all mutable limiting
          state externalised to Redis, so instances may be added or removed
          freely behind a load balancer;
        \item performing counter updates as single atomic operations, so
          coordination cost does not grow with instance count;
        \item writing request logs asynchronously off the request path, so
          that persistence latency does not enter the client's critical path;
        \item indexing log queries on (\texttt{api\_id}, \texttt{timestamp})
          and serving the dashboard from pre-aggregated summaries with
          pagination rather than raw scans.
    \end{itemize}

    \item \textbf{Operationally deployable (practically):} the system runs on
      commodity managed infrastructure:
    \begin{itemize}[leftmargin=0.7cm]
        \item managed Redis and managed PostgreSQL instances;
        \item containerised deployment of the gateway, permitting horizontal
          replication;
        \item static hosting of the React frontend on a CDN-backed platform;
        \item a CI/CD pipeline targeting staging and then production.
    \end{itemize}
\end{enumerate}

\section{Engine availability heuristic}
A mature set of ``engines'' exists for every layer of this proposition, which
is what makes the scope realistic within a course timeline:
\begin{itemize}[leftmargin=0.7cm]
    \item Spring Boot and Spring Cloud Gateway for the reverse-proxy and REST
      layers.
    \item Bucket4j as a maintained implementation of the token-bucket and
      window algorithms, with an existing Redis integration.
    \item Redis for atomic, low-latency counter operations.
    \item PostgreSQL with Spring Data JPA for durable relational storage.
    \item Spring Security with JWT for token-based authentication.
    \item React with Recharts for the dashboard and its visualisations.
    \item JUnit and Testcontainers for integration tests against real Redis and
      PostgreSQL instances.
    \item A CI runner (GitHub Actions) and free-tier staging targets.
\end{itemize}
We use these mature components deliberately, so that our engineering effort is
spent on correct distributed enforcement, instrumentation and measurement,
rather than on reimplementing infrastructure.

\section{Project scope and deliverables %
\BvrHint{\ldots iteration-friendly}}
\subsection{Initial deliverable %
\BvrHint{\ldots first iteration}}
\begin{itemize}[leftmargin=0.7cm]
    \item Working reverse proxy: request forwarded to a configured target API
      and response relayed.
    \item Token-bucket enforcement backed by Redis, returning HTTP~429 with
      \texttt{Retry-After} on rejection.
    \item User signup and login with JWT; API registration with a configurable
      limit.
    \item Request logging to PostgreSQL with timestamps and outcome, sufficient
      to compute the primary metric.
    \item CI: build, backend unit tests, and linting on every push.
    \item CD to staging as a single pipeline job.
\end{itemize}

\subsection{Subsequent deliverables}
\begin{itemize}[leftmargin=0.7cm]
    \item Sliding-window algorithm as a selectable policy per API.
    \item React analytics dashboard: allowed versus blocked counts, traffic
      over time, per-API breakdown.
    \item Request log view with filtering and pagination.
    \item Multi-instance validation run demonstrating enforcement accuracy
      across two gateways.
    \item \emph{(Stretch)} Request replay: re-issue a specific blocked or
      failed request from the log.
    \item \emph{(Stretch)} Live side-by-side comparison of the two algorithms
      under identical load.
\end{itemize}

\section{Risks and mitigations}
\begin{itemize}[leftmargin=0.7cm]
    \item \textbf{Counter store as a single point of failure:} Redis
      unavailability would disable enforcement. \emph{Mitigation:} define and
      implement an explicit fail-open/fail-closed policy, surface the degraded
      state in the dashboard, and test the failure path in CI.
    \item \textbf{Latency budget:} an additional network hop plus a counter
      lookup could exceed the acceptable overhead. \emph{Mitigation:}
      co-locate gateway and Redis in the same region, keep logging off the
      request path, and treat p95 overhead as a build-gating metric.
    \item \textbf{Race conditions in distributed counting:} the defining
      correctness risk of the project. \emph{Mitigation:} rely on atomic
      server-side operations rather than read-modify-write, and make
      multi-instance concurrent load tests part of the automated suite from
      iteration one.
    \item \textbf{Credential and secret handling:} JWT signing keys and
      registered-API credentials are sensitive. \emph{Mitigation:} store
      secrets in platform-managed environment configuration, never in the
      repository; store only what is needed and hash user credentials.
    \item \textbf{Proxy abuse:} an open proxy can be misused to reach arbitrary
      hosts. \emph{Mitigation:} permit forwarding only to targets registered by
      an authenticated owner, and validate target URLs on registration.
    \item \textbf{Scope creep from stretch features:} replay and live
      comparison are attractive but non-essential. \emph{Mitigation:} they are
      explicitly scheduled last and gated on the primary metric being met.
    \item \textbf{Free-tier deployment limits:} cold starts and connection
      caps on free tiers can distort latency measurements.
      \emph{Mitigation:} warm the service before measurement runs and report
      the deployment tier alongside every figure.
\end{itemize}

\section{Summary}
This project proposes API Gateway, a deployable, self-service rate-limiting
gateway that lets a developer place enforceable traffic control in front of an
existing API without modifying it. It meets the course criteria by
emphasising time-to-value---a useful proxy-plus-limiter is releasable in the
first iteration---by using CI/CD to support frequent safe releases, and by
defining evaluation in measurable, attributable terms, principally enforcement
accuracy under concurrent load with added latency as a bounding constraint.
It relies throughout on mature components and published algorithms, keeping
the engineering contribution in correct distributed implementation,
instrumentation and scalability readiness rather than in research novelty.

\end{document}